% Preámbulo común de los manuales de TRAZAFI (Titulación de licenciatura y Posgrado).
% Lo carga el preambulo.tex de cada nivel, que antes define \nivelmanual y la ruta de sus imágenes.
\usepackage[spanish,provide=*]{babel}
\usepackage[T1]{fontenc}
\usepackage[utf8]{inputenc}
\usepackage{lmodern}
\usepackage[scaled=0.95]{helvet}
\renewcommand{\familydefault}{\sfdefault}
\usepackage[a4paper,margin=2.2cm,headheight=14pt]{geometry}
\usepackage{graphicx}
\usepackage{xcolor}
\usepackage{longtable}
\usepackage{booktabs}
\usepackage{array}
\usepackage{enumitem}
\usepackage{fancyhdr}
\usepackage{float}
\usepackage{multicol}
\usepackage{caption}
\usepackage[most]{tcolorbox}
\usepackage[hidelinks,pdfusetitle]{hyperref}
\definecolor{uaemex}{HTML}{1F5E3B}
\definecolor{oro}{HTML}{B8860B}
\definecolor{alerta}{HTML}{9B2C2C}
\definecolor{gris}{HTML}{5F6B66}
\setlength{\parindent}{0pt}
\setlength{\parskip}{0.55em}
\setlength{\emergencystretch}{2em}
\setlist{itemsep=2pt,topsep=3pt}
\captionsetup{font=small,labelfont={bf,color=uaemex},labelsep=period,justification=centering}
\pagestyle{fancy}
\fancyhf{}
\fancyhead[L]{\small\color{gris}\tituloCorto}
\fancyhead[R]{\small\color{gris}\nivelmanual}
\fancyfoot[C]{\small\thepage}
\renewcommand{\headrulewidth}{0.4pt}
\newcommand{\tituloCorto}{Manual}
\newcommand{\fechaversion}{Noviembre de 2026}

% ---- Portada: \portada{Título}{Actor}{Modalidad u opción}
\newcommand{\portada}[3]{%
  \renewcommand{\tituloCorto}{#1}%
  \hypersetup{pdftitle={#1 · #3},pdfauthor={Facultad de Ingeniería, UAEMéx}}%
  \begin{titlepage}
    \centering
    {\large\color{uaemex}\textbf{Universidad Autónoma del Estado de México}\par}
    \vspace{2pt}{\large\color{uaemex}Facultad de Ingeniería\par}
    \vspace{2.2cm}
    \includegraphics[width=3.2cm]{logo_traza}\par
    \vspace{0.6cm}{\small\color{gris}Plataforma Traza FI\par}
    \vspace{1.6cm}
    {\Huge\bfseries\color{uaemex}#1\par}
    \vspace{0.7cm}
    {\LARGE #3\par}
    \vspace{0.4cm}
    {\large\color{gris}\nivelmanual\par}
    \vfill
    \begin{tabular}{rl}
      \textbf{Para:} & #2\\
      \textbf{Versión:} & \fechaversion\\
    \end{tabular}\par
    \vspace{1cm}
    {\footnotesize\color{gris}Las pantallas son reales; los nombres, cuentas y títulos que aparecen en ellas son ficticios\\y provienen de una base de pruebas.\par}
  \end{titlepage}
  \setcounter{page}{2}}

% ---- Capturas de pantalla
\newcommand{\captura}[2]{%
  \begin{figure}[H]\centering
    \fcolorbox{gris!40}{white}{\includegraphics[width=0.94\linewidth,height=0.56\textheight,keepaspectratio]{#1}}
    \caption{#2}
  \end{figure}}
% Captura de celular (390 px de ancho): más angosta.
\newcommand{\capturamovil}[2]{%
  \begin{figure}[H]\centering
    \fcolorbox{gris!40}{white}{\includegraphics[width=0.36\linewidth,height=0.6\textheight,keepaspectratio]{#1}}
    \caption{#2}
  \end{figure}}
% Dos capturas de celular lado a lado.
\newcommand{\capturasmovil}[4]{%
  \begin{figure}[H]\centering
    \begin{minipage}[t]{0.46\linewidth}\centering
      \fcolorbox{gris!40}{white}{\includegraphics[width=0.78\linewidth,height=0.55\textheight,keepaspectratio]{#1}}
      \captionof{figure}{#2}\end{minipage}\hfill
    \begin{minipage}[t]{0.46\linewidth}\centering
      \fcolorbox{gris!40}{white}{\includegraphics[width=0.78\linewidth,height=0.55\textheight,keepaspectratio]{#3}}
      \captionof{figure}{#4}\end{minipage}
  \end{figure}}

% ---- Recuadros
\newtcolorbox{tencuenta}{enhanced,breakable,colback=uaemex!5,colframe=uaemex,boxrule=0.6pt,arc=2pt,left=6pt,right=6pt,top=4pt,bottom=4pt,
  title={Ten en cuenta},fonttitle=\bfseries\small,coltitle=white,colbacktitle=uaemex,fontupper=\small}
\newtcolorbox{algosalemal}{enhanced,breakable,colback=alerta!4,colframe=alerta,boxrule=0.6pt,arc=2pt,left=6pt,right=6pt,top=4pt,bottom=4pt,
  title={Si algo sale mal},fonttitle=\bfseries\small,coltitle=white,colbacktitle=alerta,fontupper=\small}
\newtcolorbox{quevasahacer}{enhanced,colback=oro!8,colframe=oro,boxrule=0.6pt,arc=2pt,left=8pt,right=8pt,top=6pt,bottom=6pt,
  title={Qué vas a hacer},fonttitle=\bfseries,coltitle=white,colbacktitle=oro}
% Botón o texto de la interfaz, tal como aparece en pantalla.
\newcommand{\boton}[1]{\textbf{«#1»}}
